\documentclass[10pt]{article}
\usepackage[T1]{fontenc}
\usepackage[utf8]{inputenc}
\usepackage[french]{babel}
\usepackage{amsmath,amssymb,mathrsfs,mathpazo}
\usepackage{geometry}
\geometry{a4paper,margin=2cm}
\usepackage{graphicx,array,tabularx,tabularray,colortbl,multirow,multicol,booktabs,enumitem,pifont,ulem,fancybox,diagbox,cancel}
\usepackage{pstricks,pstricks-add,pst-plot,pst-tree,pst-node,pst-text,pst-func,pst-math,pst-bspline,pst-3dplot,pst-eucl,pst-solides3d}
\let\pscalculate\undefined
\usepackage{pst-bezier}
\usepackage{tikz,pgfplots,tkz-tab}
\usetikzlibrary{arrows,arrows.meta,patterns,calc,positioning,trees,decorations.markings,shapes.geometric,shapes.misc,shapes.arrows,intersections,decorations.pathmorphing,decorations.pathreplacing,backgrounds,shadows,babel}
\pgfplotsset{compat=1.18}
\usepgfplotslibrary{fillbetween}
\usepackage[np]{numprint}
\usepackage{listings,makecell,pas-tableur,twoopt}
\let\barre\undefined
\usepackage[pstricks]{tablvar}
\definecolor{darkgreen}{RGB}{0,100,0}
\providecommand{\red}{\color{red}}
\providecommand{\blue}{\color{blue}}
\providecommand{\green}{\color{green}}
\providecommand{\black}{\color{black}}
\providecommand{\white}{\color{white}}
\providecommand{\gray}{\color{gray}}
\providecommand{\cyan}{\color{cyan}}
\providecommand{\magenta}{\color{magenta}}
\providecommand{\yellow}{\color{yellow}}
\providecommand{\violet}{\color{violet}}
\providecommand{\orange}{\color{orange}}
\providecommand{\euro}{\texteuro}
\providecommand{\EUR}{\texteuro}
\providecommand{\ieme}{\textsuperscript{e}}
\providecommand{\degres}{\ensuremath{^\circ}}
\providecommand{\degree}{\ensuremath{^\circ}}
\providecommand{\vv}[1]{\overrightarrow{#1}}
\providecommand{\Coord}[1]{#1}
\providecommand{\decosix}{}
\providecommand{\starredbullet}{\ensuremath{\bullet}}
\pagestyle{empty}
\setlength{\parindent}{0pt}
\usepackage{adjustbox}\DeclareUnicodeCharacter{2236}{:}\DeclareUnicodeCharacter{2190}{\ensuremath{\leftarrow}}
\providecommand{\R}{\mathbb{R}}
\providecommand{\N}{\mathbb{N}}
\providecommand{\Z}{\mathbb{Z}}
\providecommand{\Q}{\mathbb{Q}}
\providecommand{\D}{\mathbb{D}}
\providecommand{\C}{\mathbb{C}}
\providecommand{\e}{{\mathrm{e}}}
\providecommand{\vect}[1]{\overrightarrow{#1}}
\providecommand{\vectt}[1]{\overrightarrow{\mathrm{#1}}}
\providecommand{\vvt}[1]{\overrightarrow{\mathrm{#1}}}
\providecommand{\np}[2][]{{#2~#1}}
\providecommand{\vv}[1]{{\overrightarrow{#1}}}
\renewcommand{\d}{{\mathrm{d}}}
\providecommand{\degres}{{^\circ}}
\providecommand{\degre}{{^\circ}}
\providecommand{\blue}{\color{blue}}
\providecommand{\red}{\color{red}}
\providecommand{\enskip}{\quad}
\providecommand{\llbracket}{\lbrack\!\lbrack}
\providecommand{\rrbracket}{\rbrack\!\rbrack}
\providecommand{\euro}{{\text{€}}}
\providecommand{\no}{{\text{n°}}}
\providecommand{\interval}[2]{{[#1;#2]}}
\providecommand{\intervFO}[2]{{[#1;#2[}}
\providecommand{\dotfill}{\ldots}

\providecommand{\ds}{\displaystyle}
\providecommand{\textding}[1]{#1}
\providecommand{\Oij}{\(\mathrm{O};\vec{\imath},\vec{\jmath}\)}
\providecommand{\Oijk}{\(\mathrm{O};\vec{\imath},\vec{\jmath},\vec{k}\)}
\providecommand{\Ouv}{\(\mathrm{O};\vec{u},\vec{v}\)}
\providecommand{\intervFF}[2]{\ensuremath{[#1;#2]}}
\renewcommand{\euro}{\ifmmode\text{\texteuro}\else\texteuro\fi}
\DeclareUnicodeCharacter{00B0}{\ifmmode^{\circ}\else\textdegree\fi}
\DeclareUnicodeCharacter{20AC}{\ifmmode\text{\texteuro}\else\texteuro\fi}
% Énoncés : Éducation nationale. Transcription et corrigés : APMEP, Denis Vergès et contributeurs. Découpage et classement : CoopMaths, Sébastien Lozano, Éric Elter et contributeurs. Adaptation de la présentation : Maths Explorer.
% CC BY-SA : https://coopmaths.fr/www/about/
% https://coopmaths.fr/alea/static/eam/2026/tex/eam_specialite_2026_06_polynesie_2.tex
\begin{document}
\noindent\textit{APMEP / CoopMaths. CC BY-SA. Présentation adaptée par Maths Explorer.}\par\medskip
\subsection*{Partie A}

On considère la fonction $f$ définie sur $\mathbb{R}$ par
\[f(x)=\left(\dfrac12 x + 1\right) \e^{-x}.\]

On note \mbox{$(\mathcal{C}_f)$} sa courbe représentative dans un repère orthonormé.

On admet que la fonction $f$ est dérivable sur $\mathbb{R}$ et on désigne par $f'$ sa fonction dérivée.

\begin{enumerate}
	\item Démontrer que, pour tout réel $x$,
	\[f'(x)=\left(-\dfrac12 x -\dfrac12\right) \e^{-x}.\]

	\item Étudier les variations de la fonction $f$ sur $\mathbb{R}$.

	\item L'affirmation \og{}Pour tout réel $x$, \mbox{$f(x)\leqslant 2$}\fg{} est-elle vraie ou fausse ?

	Justifier.
	\item Déterminer l'équation réduite de la tangente à \mbox{$(\mathcal{C}_f)$} au point d'abscisse 0.
\end{enumerate}
	\begin{center}
		\fcolorbox{black}{white}{
	\begin{minipage}[t]{0.8\linewidth}\vspace{0cm}
		\textsc{Aide aux calculs}
		\medskip

	\begin{tblr}{X[c]X[c]X[c]}
	\mbox{$\e^{1}\approx\np{2.7}$}&
	\mbox{$\e^{-1}\approx\np{0.4}$}&
	\mbox{$\e^{2}\approx\np{7.4}$}
	\end{tblr}
	\end{minipage}
}
	\end{center}

\subsection*{Partie B}

Soit $g$ la fonction polynôme du second degré définie par
\[g(x)=\np{0.5}x^2+bx+c\]
où $b$ et $c$ sont deux réels.

On crée la courbe représentée ci-dessous, constituée des deux éléments suivants :
\begin{itemize}[label=$\bullet$]
	\item Une partie de la courbe \mbox{$(\mathcal{C}_f)$} représentative de la fonction $f$ définie dans la partie A sur l'intervalle \mbox{$[-2;0]$}.

	\item Une partie de la courbe \mbox{$(\mathcal{C}_g)$} représentative de la fonction $g$  sur l'intervalle \mbox{$[0;3]$}.
\end{itemize}

\begin{center}
		\begin{tikzpicture}[xscale=1.3,yscale=1.3,rotate=0]
	\clip(-3.2,-0.6)rectangle(3.9,4.3);
	\draw[blue!30,thick](-3.2,-0.6)grid(3.9,4.3);
	\draw[line width=0.3mm,->, >=triangle 45](-3.2,0)--(3.9,0);
	\draw[line width=0.3mm,->, >=triangle 45](0,-0.6)--(0,4.3);
	\foreach \i in {-3,-2,-1,1,2,3,4}{\draw (\i,2pt)--(\i,-2pt) node[below ]{\np{\i}} (2pt,\i)--(-2pt,\i)node[left]{\np{\i}};}
	\draw(-2pt,-2pt) node[below left]{\np{0}};
	\draw[red](-1.5,1.5)node{\mbox{$(\mathcal{C}_f)$}};
	\draw[purple](2.6,2.2)node{\mbox{$(\mathcal{C}_g)$}};
	\node[draw,cross out, thick,minimum width=2pt, minimum height = 2pt] at (0,1){}; \draw(0,1)node[above right=1.5pt]{M};
	\draw[red, line width=0.3mm]plot[smooth,domain=-2:0]	(\x,{(\x/2+1)*exp(-\x)});
	\draw[purple, line width=0.3mm]plot[smooth,domain=0:3](\x,{\x*\x/2-\x/2+1});
\end{tikzpicture}

\end{center}
\vspace{5mm}
Le but de cette partie est de déterminer les valeurs des réels $b$ et $c$.
\begin{enumerate}
	\item Le point \mbox{M~$(0;1)$} appartient aux deux courbes  \mbox{$(\mathcal{C}_f)$} et  \mbox{$(\mathcal{C}_g)$}.

	Démontrer que \mbox{$c=1$}.
	\item
 Au point  \mbox{M~$(0;1)$} les courbes  \mbox{$(\mathcal{C}_f)$} et  \mbox{$(\mathcal{C}_g)$} admettent la même tangente.

 Déterminer la valeur de $b$.
\end{enumerate}

\end{document}
